\documentclass[12pt, letterpaper]{report}


\usepackage[utf8]{inputenc}
\usepackage[T1]{fontenc}
\usepackage[spanish,es-noquoting]{babel}

\usepackage[left=4cm,right=2cm,top=3cm,bottom=3cm]{geometry}
\usepackage{setspace}
\onehalfspacing

\usepackage{mathptmx}

\setlength{\parindent}{0pt}
\setlength{\parskip}{1em}

\usepackage{booktabs}
\usepackage{longtable}
\usepackage{tabularx}
\usepackage{array}
\usepackage{graphicx}
\usepackage{float}
\usepackage{caption}

\captionsetup{font=small,labelfont=bf}

\usepackage{cite}
\usepackage[hidelinks]{hyperref}
\usepackage{titlesec}

\titleformat{\chapter}[block]
  {\normalfont\bfseries\Large\centering}
  {\thechapter.}
  {0.5em}
  {\MakeUppercase}

\titleformat{\section}
  {\normalfont\bfseries\large}
  {\thesection}
  {0.5em}
  {\MakeUppercase}

\titleformat{\subsection}
  {\normalfont\bfseries\normalsize}
  {\thesubsection}
  {0.5em}
  {}

\usepackage{fancyhdr}
\pagestyle{plain}

\begin{document}

\begin{titlepage}
\centering
\setstretch{1.15}

{\MakeUppercase{Universidad Pedagógica y Tecnológica de Colombia}}\\[0.2cm]
{\MakeUppercase{Facultad de Ingeniería}}\\[0.2cm]
{\MakeUppercase{Ingeniería de Sistemas y Computación}}\\[2.5cm]

\textbf{\Large EFECTO DE UN ASISTENTE DE INTELIGENCIA ARTIFICIAL GENERATIVA EN EL TIEMPO DE DEPURACIÓN DE ERRORES DE SINTAXIS EN ESTUDIANTES NOVELES DE PROGRAMACIÓN}\\[0.5cm]
\textit{\small Título provisional — se ajusta al formular la propuesta (Fase 3)}\\[2.5cm]

\MakeUppercase{Autores}\\[0.4cm]
Barrera Romero Juan Diego\\[0.2cm]
Castro Gómez Lui Miguel\\[0.2cm]
Corredor Silva Santiago\\[0.2cm]
Galeano Tirado Juan Emmanuel

\vfill
Documento elaborado en el marco del curso Metodología de la Investigación\\[0.3cm]
Tunja, Colombia · 2026
\end{titlepage}

\newpage
\tableofcontents
\newpage
\listoftables
\newpage
\listoffigures
\newpage


\chapter{Planteamiento del problema}
\section{Antecedentes}
adhajkdsa dajkshdkajsdk dhkaas
dgdfgd
f
\section{Diagnóstico}

\section{Pronóstico}

\section{Control del pronóstico}
\chapter{Justificación} 
\chapter{Objetivos}
\chapter{Alcance de la investigación}
\chapter{Formulación de hipótesis}
\chapter{Variables}
\chapter{Marco referencial}
\chapter{Diseño metodológico} 
\chapter{Resultados y discusión}
\chapter{Conclusiones y recomendaciones}


\bibliographystyle{texstyles/IEEEtran}
\nocite{*}
\bibliography{referencias}

\end{document}